\section{Introduction}\label{sec:intro}

Straight-line distance is the default abstraction of applied spatial analysis. It is a useful simplification,
but people do not move across a featureless plane: they walk along streets, around blocks, over bridges and up
stairways, and the configuration of that network conditions whom they can easily meet. In segregation research
the abstraction enters through the definition of the ``local environment'' that spatial indices compare with the
city as a whole \citep{reardon2004measures}. Most applications define that environment with Euclidean distance.
\citet{knaap2024segregated} showed, for every metropolitan area in the United States, that replacing Euclidean
distance with walking distance along the street network raises the spatial information theory index in almost
every metro, by about a fifth on average, and reshuffles the ranking of the most segregated places. They also
showed that the size of the gap is associated with the topology of the street network.

This paper asks whether the same holds in Brazil, where there are strong reasons to expect the street network to
matter even more. Brazilian cities combine colonial cores, modernist plans, sprawling peripheral \emph{loteamentos},
self-built favelas on steep hillsides and, since the 1990s, gated condominiums (\emph{condom\'inios fechados}) that
interrupt the public street grid \citep{coy2006gated}. In the most comprehensive configurational comparison to
date, \citet{medeiros2013urbis} found that the street networks of the 27 Brazilian state capitals are among the
most fragmented of 164 cities worldwide: a ``patchwork'' of locally coherent pieces that are poorly articulated
with one another. If segregation is, in part, a restriction on the possibility of interaction
\citep{netto2015segregated}, then a patchwork city should look more segregated when the local environment is
measured along the network than when it is measured as the crow flies.

We compute racial segregation for the \nCities{} Brazilian municipalities with at least 300,000 inhabitants,
using the 2022 Census colour-or-race counts for urban census tracts and pedestrian networks from OpenStreetMap.
For each city we build two local environments that differ \emph{only} in the distance metric: the same 2-km
triangular kernel is applied to Euclidean and to network distances. We ask two questions. First, how different
are segregation levels and rankings under the two metrics (RQ1)? Second, which properties of the street network
are associated with the gap (RQ2)? To RQ2 we add a Brazil-specific predictor, a metric analogue of the
space-syntax \emph{synergy} measure that underlies Medeiros's diagnosis of fragmentation. Because network
environments contain fewer people than Euclidean environments of the same radius, part of any gap could be a
scale effect rather than a barrier effect \citep{saporito2026reconsidering}; we therefore also report a
population-matched comparison in which each tract's Euclidean environment is shrunk until it holds as many
people as its network environment.

The paper makes three contributions: the first nationwide Euclidean-versus-network comparison of segregation for
a country of the Global South; a test of whether configurational fragmentation, a long-standing theme of
Brazilian urban morphology, helps explain the gap; and a fully reproducible workflow built on the open-source
PySAL \texttt{segregation} package \citep{cortes2026segregation,cortes2020open,rey2021comparative,rey2022pysal}.
Consistent with \citet{knaap2024segregated}, we treat the comparison descriptively; a formal inferential test of
the gap is left for future work.

\section{Background}\label{sec:background}

\subsection{Measuring segregation in space}

Spatial segregation indices replace the question ``who lives in the same tract?'' with ``who lives nearby?''.
\citet{white1983measurement} introduced distance decay into the measurement of segregation, and
\citet{wong1993spatial} showed that indices based on the adjacency of areal units depend on unit size and shape.
\citet{reardon2004measures} provided a general framework in which any theoretically motivated proximity function
defines each person's local environment; among the resulting measures, the spatial information theory index $H$
has the most desirable properties. \citet{reardon2008geographic} used the bandwidth of the proximity function to
describe the geographic scale of segregation. In Brazil, \citet{feitosa2007global} developed kernel-based global
and local spatial indices and applied them to S\~ao Jos\'e dos Campos. In all of these, proximity is almost always
Euclidean.

\subsection{Street networks and social interaction}

A second literature argues that street networks structure the opportunities for contact.
\citet{grannis1998importance} showed that residential (``trivial'') streets bound neighbourhood networks, and
\citet{grannis2005tcommunities} that pedestrian-connected ``T-communities'' predict racial composition in three US
cities. \citet{roberto2018spatial} incorporated road connectivity into segregation measurement and found higher
segregation in Pittsburgh under network distance. \citet{knaap2024segregated} scaled this comparison to all US
metros and related the gap to network topology measured with OSMnx \citep{boeing2017osmnx} and momepy
\citep{fleischmann2019momepy}: circuity widens the gap, intersection density narrows it, and cyclomatic
complexity (redundancy) dampens the effect of circuity. Global comparisons show that street-network form varies
far more outside North America: grids are the exception, and circuity and orientation entropy are higher
\citep{boeing2019urban,boeing2022street}. Recently, \citet{saporito2026reconsidering} cautioned that network-based
environments of a given radius contain fewer people, so that part of the network--Euclidean difference may reflect
the population size of local environments rather than barriers to movement.

\subsection{The configuration of Brazilian street networks}

Brazilian research on urban form has long emphasised configuration. Drawing on space syntax,
\citet{holanda2002espaco} interpreted the plan of Bras\'ilia as a ``space of exception'' whose dispersed,
hierarchical layout separates social groups. \citet{medeiros2013urbis} compared axial maps of 164 cities, including
all 27 Brazilian capitals, and concluded that Brazilian cities are exceptionally fragmented: patchworks of
internally coherent pieces whose local and global structures are poorly articulated (low \emph{synergy}), a
labyrinthine configuration that hampers movement between parts of the city. \citet{netto2015segregated} reframed
segregation as a restriction on interaction produced jointly by residential location, the street network and
mobility, with evidence from a Brazilian city. Two recent studies bring the street grid directly into segregation
measurement: \citet{saboya2025chasm} propose CHASM, an exposure measure computed on street segments, and
\citet{goncalves2024uncovering} use a network-based potential-accessibility index to reveal racial and income
segregation patterns in Pelotas. At the scale of informal settlements, \citet{carvalho2023segregation} show that
favelas are themselves internally segregated. Using OSM, \citet{spadon2018topological} characterised the street
networks of the 645 municipalities of S\~ao Paulo state by topological features. Finally, the spread of gated
condominiums has privatised parts of the street network and fragmented the urban fabric \citep{coy2006gated}.
Together these works motivate three expectations: (H1) segregation is higher under network than under Euclidean
distance; (H2) the gap is larger where circuity, dead-end streets and configurational fragmentation are higher;
and (H3) the gap is smaller where intersection density and meshedness are higher.

\subsection{Race, space and access in Brazilian cities}

Racial residential segregation in Brazil is moderate by US standards but real and entangled with class
\citep{telles2004race}; a nationwide analysis of 2010 census tracts found it highest in the South and Southeast
\citep{sousafilho2023segregacao}. The consequences of location are mediated by the transport network. The IPEA
\emph{Acesso a Oportunidades} project shows that white and high-income residents of the 20 largest Brazilian
cities reach more jobs, schools and health services than black and low-income residents by every transport mode
\citep{pereira2019desigualdades,tomasiello2024racial}, and that transport policy can widen these gaps
\citep{pereira2019distributional}. \citet{bittencourt2021cumulative} show that accessibility and segregation
reinforce each other in four metropolises, and \citet{boisjoly2020accessibility} document large accessibility
differences across four cities. This literature measures travel along networks
\citep{pereira2021r5r} but takes residential segregation as given; our study brings the network into the
measurement of segregation itself.
